\SetPageTheme{story}
\PageHeader
  {Operatori aritmetici}
  {Scorul, viețile și energia cresc sau scad prin pași mici, repetați corect la fiecare acțiune.}
  {Operatori 1 din 5}

\begin{missionbox}[Narațiunea paginii]
În jocuri, cele mai dese transformări sunt adunările și scăderile: +1 punct la colectare, -1 viață la coliziune, +10 bonus la final de nivel, -5 energie la folosirea unei abilități.

Când înțelegi aceste transformări matematice, începi să vezi codul nu ca pe o simplă listă de comenzi, ci ca pe o secvență de actualizări controlate, esențiale pentru dinamica jocului.
\end{missionbox}

\begin{conceptbox}[Reținem!]
\begin{itemize}
\item \texttt{scor <- scor + 1} înseamnă un pas înainte în progresul tău.
\item \texttt{vieti <- vieti - 1} înseamnă un pas înapoi în siguranța ta.
\item \texttt{energie <- energie + k} înseamnă un salt controlat de resurse.
\end{itemize}
\end{conceptbox}

\BlocklyExperiment{cq-i1-operator-score}{Scor incremental}{Testează blocurile care cresc sau scad scorul în funcție de evenimentele din joc. Vezi cum influențează fiecare alegere starea finală.}

\SyntaxCheck{operatori}

\clearpage

\SetPageTheme{guided}
\PageHeader
  {Exerciții ghidate}
  {Începem cu exemple scurte pentru a verifica reflexele, apoi extindem spre mini-scenarii de joc.}
  {Operatori 2 din 5}

\begin{guidebox}[Nivel ghidat: Calcule rapide]
Completează valorile lipsă:
\begin{enumerate}
\item \texttt{scor <- 8}; după \texttt{scor <- scor + 2}, \texttt{scor = } \AnswerSpace{1}
\item \texttt{vieti <- 4}; după \texttt{vieti <- vieti - 1}, \texttt{vieti = } \AnswerSpace{1}
\item \texttt{monede <- 3}; după \texttt{monede <- monede + 5}, \texttt{monede = } \AnswerSpace{1}
\end{enumerate}
\AnswerSpace{2}
\end{guidebox}

\begin{guidebox}[Nivel ghidat: Traseul stării]
Pornești cu \texttt{energie <- 10}. Se execută următoarea secvență:
\begin{lstlisting}[style=codequest]
energie <- energie - 2
energie <- energie + 1
energie <- energie - 3
\end{lstlisting}
Care este valoarea energiei la final?
\AnswerSpace{4}
\end{guidebox}

\clearpage

\SetPageTheme{guided}
\PageHeader
  {Scrii transformările}
  {Acum completezi singur liniile lipsă, dar păstrezi același fir logic de actualizare.}
  {Operatori 3 din 5}

\begin{semibox}[Completează codul]
Vrem să obținem exact următoarea evoluție a scorului: \texttt{5 -> 7 -> 6 -> 10}.
\begin{lstlisting}[style=codequest]
scor <- 5
scor <- ................
scor <- ................
scor <- ................
\end{lstlisting}
\AnswerSpace{5}
\end{semibox}

\begin{semibox}[Scrie regula]
Explică în 2-3 rânduri de ce \texttt{x <- x + 1} este considerată o \textit{actualizare a stării} și nu o simplă comparație matematică de egalitate.
\AnswerSpace{4}
\end{semibox}

\SimulatorCard{cq-i1-operator-score}{Operatori și punctaj}{Relația dintre evenimentele din joc și modificările numerice constante din variabile.}{Schimbă valorile bonusului și penalizării în simulator, apoi compară două strategii de punctaj: una conservatoare și una agresivă. Care ajunge la target mai repede?}

\clearpage

\SetPageTheme{challenge}
\PageHeader
  {Rezolvi scenarii complete}
  {Rezolvi scenarii complete, cu explicații scurte și rezultate verificabile prin logică pură.}
  {Operatori 4 din 5}

\begin{solobox}[Set principal de provocări]
\begin{enumerate}
\item Pornind cu \texttt{scor <- 0}, scrie o secvență de 6 linii prin care să ajungi exact la valoarea 20.
\item Pornind cu \texttt{vieti <- 5}, scrie o secvență de 4 linii prin care să ajungi la valoarea 1.
\item Pornind cu \texttt{energie <- 12}, creează o secvență de actualizări care să se termine exact la valoarea 9.
\end{enumerate}
\AnswerSpace{10}
\end{solobox}

\clearpage

\SetPageTheme{challenge}
\PageHeader
  {Explici de ce funcționează}
  {Un programator bun nu doar calculează corect. El știe să explice clar de ce acel calcul este logic.}
  {Operatori 5 din 5}

\begin{solobox}[Problema jurnalului de bord]
Tabela de scoruri înregistrată într-o sesiune de test este: \texttt{2, 6, 5, 11, 8}. Scrie o secvență de operații care poate produce exact această evoluție a scorului.
\AnswerSpace{8}
\end{solobox}

\begin{solobox}[Reflecție]
Ce greșeli apar cel mai des atunci când lucrezi cu operații de tipul \texttt{variabilă <- variabilă + ceva}? De ce este ușor să pierzi urma valorii corecte?
\AnswerSpace{6}
\end{solobox}

\clearpage
